\begin{figure}[t]
\centering
\includegraphics[width=\linewidth]{static/images/teaser_tex.pdf}

\vspace{-3pt}
\caption{\textbf{Typicality bias across the modern vision stack.}
\emph{Left:} Three \emph{anchor classes} (target object categories we examine): \textit{\textcolor{anchorcolor}{skateboard}}, \textit{\textcolor{anchorcolor}{dog}}, and \textit{\textcolor{anchorcolor}{cow}}, shown in their {\color{typcolor}\textbf{typical}} company (top row: with a \textit{\textcolor{typcolor}{person}}, with a \textit{\textcolor{typcolor}{frisbee}}, with a \textit{\textcolor{typcolor}{bird}}) and in {\color{atypcolor}\textbf{atypical}} surroundings (bottom row: with a \textit{\textcolor{atypcolor}{cat}}, with a \textit{\textcolor{atypcolor}{spoon}}, with a \textit{\textcolor{atypcolor}{clock}}). \emph{Right:} Each model's score on the anchor classes in COCO \texttt{val} images where they appear with typical vs.\ atypical partners (NPMI quintile splits, \S\ref{sec:diagnostic}). The performance gap $\textcolor{atypcolor}{\Delta} =$ (\textcolor{black}{Typical} -- \textcolor{black}{Atypical}) persists across all five model families and three tasks we examine, ranging from $\textcolor{atypcolor}{\Delta = 5.4}$~p.p.\@ (Co-DINO Swin-L) to $\textcolor{atypcolor}{\Delta = 11.0}$~p.p.\@ (OpenCLIP).}
\label{fig:teaser}
\vspace{-12pt}
\end{figure}
